%!TEX program = lualatex
% Arquivo: Template_Estat_Texto.tex

\documentclass[a4paper,12pt]{article}

% =========================================================
% Idioma e tipografia (LuaLaTeX)
% =========================================================
\usepackage[portuguese]{babel}
\usepackage{fontspec}
\setmainfont{Latin Modern Roman}

% =========================================================
% Layout
% =========================================================
\usepackage[top=3cm,bottom=2.5cm,left=3cm,right=2.5cm]{geometry}
\usepackage{setspace}
\onehalfspacing
\usepackage{indentfirst}
\setlength{\parindent}{1.25cm}

% =========================================================
% Matemática (align*, etc.)
% =========================================================
\usepackage{amsmath}
\usepackage{amssymb}

% =========================================================
% Tabelas, listas, figuras
% =========================================================
\usepackage{booktabs}
\usepackage{array}
\usepackage{enumitem}
\setlist[itemize]{noitemsep,topsep=3pt}
\setlist[enumerate]{noitemsep,topsep=3pt}

\usepackage{graphicx}
\usepackage{float}
\usepackage{pdflscape}
\graphicspath{{./}{./imagens/}{./figuras/}}

% =========================================================
% Links
% =========================================================
\usepackage{hyperref}
\usepackage{xurl}
\hypersetup{
  colorlinks=true,
  linkcolor=black,
  urlcolor=blue,
  citecolor=black,
  pdfauthor={},
  pdftitle={}
}

\usepackage{siunitx}

% =========================================================
% Texto de preenchimento (lorem)
% =========================================================
\usepackage{lipsum}

% =========================================================
% Dados da capa
% =========================================================
\newcommand{\Instituicao}{Universidade do Estado do Rio de Janeiro -- Uerj}
\newcommand{\Curso}{Química VIII / Oceanografia}
\newcommand{\Professor}{Professor: Ítalo Vieira}
\newcommand{\Aluno}{Alunos: Israel Carneiro\\
José Mauro Xavier Elsas\\
Miguel Gayani}
\newcommand{\Cidade}{Rio de Janeiro}
\newcommand{\Ano}{2026}

\title{\textbf{Título do Documento}}
\author{\Aluno}
\date{\Cidade --- \Ano}

\begin{document}

% =========================================================
% Capa (página de rosto)
% =========================================================
\begin{titlepage}
\centering
{\large \Instituicao \par}
\vspace{0.5cm}
{\large \Curso \par}
\vspace{1.0cm}

{\LARGE \textbf{Prática 3} \par}
\vspace{0.6cm}
{\Large PREPARO E PADRONIZAÇÃO DE SOLUÇÃO\\ DE HCl 0,1 mol L-1 \par}

\begin{figure}[H]
    \centering
    \includegraphics[width=0.3\linewidth]{./uerj.jpg}
\end{figure}

\vfill

\begin{flushright}
\Aluno\\
\Professor
\end{flushright}

\vspace{1.0cm}
{\Cidade --- \Ano \par}
\end{titlepage}

% =========================================================
% Sumário
% =========================================================
% \tableofcontents
% \newpage

\section{INTRODUÇÃO}

Os reagentes mais comumente usados para o preparo de soluções padrão de ácidos são o ácido
clorídrico, o sulfúrico e o perclórico. A solução de ácido clorídrico é a preferida, pois:

\begin{itemize}
    \item é estável indefinidamente;
    \item pode ser usada em presença da maior parte dos cátions sem sofrer interferência devido à formação de sais pouco solúvel.
\end{itemize}


A solução de HCl concentrada tem, aproximadamente 37\% (m/m) e densidade \SI{1,19}{g/mL}. Esta
solução não é um padrão primário. A solução de HCl \SI{0,1}{mol\cdot L^{-1}} é, portanto, preparada pela diluição de um volume conhecido do ácido concentrado.

\section{Material Utilizado}
\begin{enumerate}[label=\alph*.]
    \item Cilindro de 500mç
    \item Balão volumétrico de 250ml
    \item Becher de 200  300ml
    \item Bastão de vidro
    \item Funil
    \item 2 Frascos de vidro ou plástico com capacidade de 500 mL;
\end{enumerate}

\section{Reagentes}

\begin{enumerate}[label=\alph*.]
\item Ácido clorídrico (HCl) concentrado p.a..
\item Solução do indicador alaranjado de metila.
\item Solução de tetraborato de sódio deca-hidratado, (Na2B4O7.10H2O) bórax, p. a.
\end{enumerate}

\section{Procedimento Experimental}



% =========================================================
% Referências (entra no sumário)
% =========================================================
\newpage
\phantomsection
\addcontentsline{toc}{section}{Referências}
\begin{thebibliography}{99}

\bibitem{uerj_quimica8_2017}
LAGO, Dalva Cristina Baptista do;
GOIS, Jefferson Santos de;
LANGONE, Marta Antunes Pereira;
SENNA, Lilian Ferreira de.
\textbf{Química VIII — Oceanografia: aulas práticas}.
Rio de Janeiro: Universidade do Estado do Rio de Janeiro, Instituto de Química, Departamento de Química Analítica, 2017.

\end{thebibliography}

\end{document}
